% Augur --- NSDI/SOSP submission draft (acmart sigconf).
% Compiles standalone in "nonacm" draft mode (Overleaf or a TeX Live install
% with texlive-publishers, which ships acmart + ACM-Reference-Format.bst).
% Requires the shared bibliography pmcsched-refs.bib in the same directory,
% and the pgfgantt and array packages (both ship with TeX Live).
% For a double-blind submission, replace the \documentclass line with:
%   \documentclass[sigconf,review,anonymous]{acmart}
% and fill in the \acmConference / \setcopyright metadata.

\documentclass[sigconf,nonacm]{acmart}
\settopmatter{printacmref=false}
\setcopyright{none}
\renewcommand\footnotetextcopyrightpermission[1]{}

\usepackage{booktabs}
\usepackage{array}
\usepackage{pgfgantt}

\begin{document}

\title{Augur: Predictive Counter-Driven Scheduling for Tiered and
CXL-Attached Memory}

\author{Anonymous Author(s)}
\affiliation{%
  \institution{Affiliation withheld for review}
  \country{}
}

\begin{abstract}
Interference-aware schedulers react to contention after it appears. This is
acceptable when the corrective action is cheap---reallocating a core takes
microseconds---but it breaks down on tiered and CXL-attached memory, where the
contention that matters is link and controller bandwidth and the corrective
actions (page migration across tiers, cross-socket rebalancing) cost orders of
magnitude more. By the time a reactive controller observes the contention, the
tail latency it was protecting has already been spent. We argue that hardware
performance-monitoring-counter (PMC) \emph{time series} carry enough leading
signal to predict phase transitions before they materialize, and we present
Augur, a scheduler that forecasts contention from the counter stream and
\emph{pre-positions} both threads and pages across the memory hierarchy. Augur
co-designs placement with tiering: rather than letting an independent page-mover
chase hot pages after the fact, the scheduler anticipates a phase and stages the
working set on the right tier ahead of demand. We target memory-bound
microservices on real CXL hardware and evaluate against reactive interference
control and transparent tiering, with the central question being whether
prediction beats reaction once corrective actions are expensive.
\end{abstract}

\maketitle

\section{Introduction}
Datacenter resource managers have driven the per-request correction loop down to
microseconds. Shenango~\cite{ousterhout2019shenango} reallocates cores at
microsecond granularity to keep latency-sensitive services responsive, and
Caladan~\cite{fried2020caladan} extends this to interference, reacting to
memory-bandwidth and hyperthread contention by moving cores almost as fast as the
contention arises. The implicit assumption is that the corrective action is cheap
relative to the contention's duration.

Tiered and CXL-attached memory breaks that assumption. CXL Type-3 memory adds a
second, slower, lower-bandwidth tier behind a load-store
interface~\cite{cxl-spec}, and the systems response has been transparent page
placement and pooling that migrate hot data toward fast
memory~\cite{maruf2023tpp,li2023pond}. But migration is expensive and
asymmetric: moving a working set across the tier boundary costs far more than
stealing a core, and the bandwidth contention that triggers it lives on a link
the scheduler does not currently observe. A reactive controller therefore pays
twice---once for the latency incurred before it notices, and again for the
high-cost migration it must then perform---and it can thrash if phases alternate
faster than migrations complete.

Our thesis is that the PMU is a \emph{leading} indicator, not merely a lagging
one: the counter time series preceding a phase transition (rising LLC miss rate,
shifting bandwidth split across tiers, changing stall composition) is predictive
enough to act on \emph{before} the expensive contention materializes. Augur
forecasts the transition and pre-positions threads and pages so that the
corrective action is already complete when demand arrives.

\paragraph{Contributions.}
\begin{enumerate}
\item Evidence that counter time series predict memory-tier phase transitions far
enough ahead to amortize migration cost---the precondition for the whole design.
\item A lightweight online predictor over the PMC stream and a scheduler that
turns its forecasts into joint thread-placement and page-staging decisions.
\item A co-design of scheduling and tiering in which the scheduler, not an
independent page-mover, owns the anticipatory placement decision.
\item An evaluation on real CXL hardware isolating the benefit of prediction over
reaction.
\end{enumerate}

\section{Background and Motivation}
\paragraph{Why reaction is too slow here.}
On a monolithic DRAM system, the cost of mispredicting contention is a few
microseconds of recovery. On a tiered system, a misplaced working set incurs the
full fast-vs-slow latency and bandwidth gap on every access until a migration
completes, and the migration itself consumes bandwidth that competes with the
very workload it is trying to help. The economics of correction have changed, and
with them the case for prediction.

\paragraph{Counters as leading indicators.}
The contention-aware scheduling literature used counters as an instantaneous
classifier of a co-runner's badness~\cite{zhuravlev2010contention}, and even the
earliest counter-driven scheduling work sampled behavior to decide
placement~\cite{snavely2000symbiotic,suh2002memory}. Augur's departure is
temporal: it models the counter \emph{trajectory}. Many datacenter services are
strongly phased at the request and epoch level, which is precisely the structure
a short-horizon predictor can exploit.

\paragraph{Planned motivating measurement.}
We will instrument a memory-bound microservice on a CXL testbed and show (i) that
tail-latency spikes coincide with bandwidth-split phase transitions, and (ii)
that those transitions are preceded by a detectable counter signature with enough
lead time to stage pages before the spike---quantifying the lead-time budget the
predictor must hit to be worthwhile.

\section{Design}
\paragraph{Predictor.}
Augur runs a lightweight online model over a rotating set of PMCs (per-tier
bandwidth via RDT MBM~\cite{intel-rdt}, LLC miss rate, IPC, and stall
composition). The model is deliberately small---change-point detection over the
counter stream plus a short-horizon classifier of the next phase---so that it
runs in the scheduler path with bounded overhead. The output is a forecast of the
next phase and a confidence, not a control action.

\paragraph{Anticipatory actuation.}
On a confident forecast, Augur (i) places or migrates threads to the cache/memory
domain that the predicted phase will favor, and (ii) stages the predicted hot set
on fast memory ahead of demand, demoting cold pages to make room. Crucially the
scheduler owns this decision so that placement and tiering agree; an independent
page-mover reacting to access bits would fight the scheduler's intent.

\paragraph{Guarding against misprediction.}
Prediction can be wrong, and acting on a wrong forecast wastes bandwidth. Augur
gates actuation on confidence and on a cost model: it only pre-migrates when the
expected saved tail latency exceeds the migration's bandwidth cost, and it falls
back to reactive behavior when confidence is low, so that a bad predictor never
does worse than a reactive baseline.

\paragraph{Mechanism.}
Augur is built on \texttt{sched\_ext}~\cite{linux-schedext} for placement and on
the kernel's tiering interfaces for page staging, keeping it deployable on a
stock kernel.

\section{Evaluation Plan}
\paragraph{Hardware.} A server with CXL Type-3 memory (fast local DRAM + slower
CXL tier), with a NUMA-latency-emulated configuration as a cross-check where real
CXL parts are unavailable.

\paragraph{Workloads.} Memory-bound microservices (key-value stores, a search
leaf), a recommendation-style embedding lookup, and a graph-analytics job---all
chosen for phased, bandwidth-sensitive behavior.

\paragraph{Baselines.} Caladan (reactive interference control), transparent
tiering in the style of TPP~\cite{maruf2023tpp} and pooling \`a la
Pond~\cite{li2023pond}, and stock Linux tiering. The pivotal comparison is
predict (Augur) vs.\ react (everything else) at matched migration budgets.

\paragraph{Metrics and ablations.} p99/p99.9 latency, fast-tier hit rate, wasted
migration bandwidth, and machine utilization at fixed SLO. Ablations: reaction-only
Augur (predictor disabled), prediction without the cost gate, and full Augur;
lead-time sensitivity; and behavior under deliberately mispredicted phases to show
the fallback never regresses below reactive.

\section{Related Work}
Augur builds on microsecond-scale datacenter
scheduling~\cite{ousterhout2019shenango,fried2020caladan} and on counter-driven
contention awareness~\cite{zhuravlev2010contention,snavely2000symbiotic,
suh2002memory,fedorova2007isolation}, but inverts their temporal stance from
reactive to predictive. It is complementary to transparent tiering and CXL
pooling~\cite{maruf2023tpp,li2023pond}: those systems decide \emph{where} pages
should live, while Augur decides \emph{when} to move them and the threads that use
them, driven by a forecast rather than by observed access after the fact.
Table~\ref{tab:related} places Augur against the partitioning and scheduling
lines, including its companion design Concerto.

\begin{table*}[t]
\centering
\footnotesize
\setlength{\tabcolsep}{6pt}
\renewcommand{\arraystretch}{1.3}
\begin{tabular}{@{}>{\raggedright}p{2.6cm} >{\raggedright}p{3.1cm} >{\raggedright}p{3.5cm} l >{\raggedright}p{2.3cm} >{\raggedright\arraybackslash}p{2.2cm}@{}}
\toprule
\textbf{System} & \textbf{Signal} & \textbf{Actuator(s)} & \textbf{Mode} & \textbf{Commodity path?} & \textbf{Target HW} \\
\midrule
Suh et al.\ '02~\cite{suh2002memory}        & Cache miss-rate curves        & Scheduling + partitioning (proposed) & React   & No (custom HW)   & Uniform CMP \\
Fedorova et al.\ '07~\cite{fedorova2007isolation} & IPC (cache-fair)        & Scheduling (timeslice)               & React   & Research kernel  & CMP \\
Zhuravlev et al.\ '10~\cite{zhuravlev2010contention} & LLC miss rate        & Scheduling (placement)               & React   & Research         & Multicore \\
Heracles '15~\cite{lo2015heracles}          & Latency probes, util.\        & Partitioning (CAT, DVFS, power)      & React   & Yes              & Uniform \\
PARTIES '19~\cite{chen2019parties}          & Per-service QoS               & Partitioning (multi-resource)        & React   & Yes              & Uniform \\
Caladan '20~\cite{fried2020caladan}         & Mem.\ bandwidth, req.\ time   & Scheduling (fast core realloc.)      & React   & Kernel module    & Uniform \\
Concerto (companion)                        & Full PMC + CMT/MBM            & Scheduling + partitioning (joint)    & React   & \texttt{sched\_ext} & Chiplet / CXL \\
\textbf{Augur} (this work)                  & PMC \emph{time series}        & \textbf{Scheduling + page tiering}   & \textbf{Predict} & \texttt{sched\_ext} & \textbf{Tiered / CXL} \\
\bottomrule
\end{tabular}
\caption{Where Augur sits relative to counter-driven scheduling and QoS
partitioning. Prior systems commit to a single actuator and react after
contention appears; Augur instead forecasts phase transitions from the counter
trajectory and pre-positions threads and pages, while the companion Concerto
design co-actuates scheduling and partitioning. ``Commodity path?'' indicates
whether the technique was shown in a deployable commodity-kernel path rather than
simulation or a research OS.}
\label{tab:related}
\end{table*}

\section{Research Plan and Threats to Validity}
The make-or-break threat is whether prediction genuinely beats reaction rather
than merely matching it; the evaluation must show a regime where the lead-time
budget is met and the cost-gated actuator wins, and must demonstrate graceful
fallback when the predictor is wrong. Secondary threats are predictor overfitting
to phase-regular benchmarks (addressed with irregular, adversarial phase
schedules) and the availability of representative CXL hardware (addressed with an
emulated cross-check). The 18-month plan is shown in Figure~\ref{fig:gantt}.

\begin{figure*}[t]
\centering
\resizebox{\textwidth}{!}{%
\begin{ganttchart}[
  hgrid, vgrid,
  x unit=0.82cm,
  y unit title=0.6cm,
  y unit chart=0.6cm,
  title height=1,
  bar height=0.6,
  bar/.style={fill=black!22},
  group/.style={fill=black!55},
  milestone/.style={fill=black!78},
  bar label font=\small,
  group label font=\small\bfseries,
  milestone label font=\small\itshape,
  title label font=\small
]{1}{18}
\gantttitle{Month}{18} \\
\gantttitlelist{1,...,18}{1} \\
\ganttgroup{Measurement}{1}{3} \\
\ganttbar{Counter-trajectory / lead-time study}{1}{3} \\
\ganttgroup{Predictor}{3}{8} \\
\ganttbar{Online phase predictor}{3}{7} \\
\ganttbar{Cost-gated actuation model}{6}{8} \\
\ganttgroup{Mechanism}{5}{11} \\
\ganttbar{sched\_ext placement + page staging}{5}{10} \\
\ganttbar{Tiering co-design}{8}{11} \\
\ganttgroup{Evaluation}{10}{15} \\
\ganttbar{CXL hardware + emulation}{10}{14} \\
\ganttbar{Predict-vs-react ablations}{13}{15} \\
\ganttgroup{Writing}{14}{17} \\
\ganttbar{Draft}{14}{16} \\
\ganttmilestone{Submission}{17} \\
\ganttmilestone{Artifact evaluation}{18}
\end{ganttchart}%
}
\caption{Augur 18-month research plan.}
\label{fig:gantt}
\end{figure*}

% Force the full PMC-driven-scheduling reference set to appear in the
% bibliography even where a given citation is not made in-text. Remove this
% line if you want only the works actually cited above.
\nocite{snavely2000symbiotic,snavely2002priorities,bulpin2005hyperthreading,
tam2007clustering,knauerhase2008osobservations,suh2002memory,fedorova2007isolation,
banikazemi2008pam,zhuravlev2010contention,blagodurov2010contentionaware,
merkel2010resourceconscious,blagodurov2011numa}

\bibliographystyle{ACM-Reference-Format}
\bibliography{pmcsched-refs}

\end{document}
